\begin{alist}
\item Γνωρίζουμε ότι μια συνάρτηση της μορφής $g(x) = 2\eta\mu x$ έχει ελάχιστη τιμή $-2$ και μέγιστη $2$. Άρα, η συνάρτηση $f(x) = 2\eta\mu x + 1$ έχει ελάχιστη τιμή $-2 + 1 = -1$ και μέγιστη $2 + 1 = 3$.

\item Η τιμή του $x$ για την οποία η συνάρτηση $f$ παρουσιάζει μέγιστη τιμή είναι η λύση της εξίσωσης:
\begin{align*}
f(x) = 3 &\Leftrightarrow 2\eta\mu x + 1 = 3 \Leftrightarrow \\
2\eta\mu x &= 2 \Leftrightarrow \eta\mu x = 1 \Leftrightarrow \\
\eta\mu x &= \eta\mu\frac{\pi}{2} \Leftrightarrow \\
x = \frac{\pi}{2} &\text{ ή } x = \pi - \frac{\pi}{2} = \frac{\pi}{2}.
\end{align*}
Άρα, η συνάρτηση $f$ παρουσιάζει μέγιστη τιμή για $x = \dfrac{\pi}{2}$.
\end{alist}
